\chapter{Collection of Seeds}

Not all seeds are as easy to collect as Aquilegia and Jeffersonia where the
seed capsule need only to be turned over and the seeds pour out. Seeds of
Penstemon are enclosed ma hard capsule with a sharp beak; The procedure with
these is to lightly crush the capsules with pliers after which the seed fails out with some
mechanical manipulation. Separation of .the seed from the capsule debris is
accomplished by shaking the mixture lightly in a cupped sheet of smooth paper. The
debris rises to the top and the seed settles to the bottom. The sheet of paper is laid flat
and the debris brushed to one side and off the paper using a half inch wide soft brush.
The process is repeated several times. This brushing procedure has been much more
effective with small samples of seeds than any version of blowing or the ancient art of
threshing and winnowing.

Seeds of Campanulaceae and Lobelia are shed from pores in the sides of the
seed capsule. Shaking the seed capsules vigorously in a paper bag will form a pool
of seed at the bottom. These can be given a final cleaning by the brushing technique
described above. Seeds of Acantholimon and Asteraceae are easy enough to collect,
but there is much voluminous chaff and empty seed coats. These are not easily
removed in the dry state, but after placing on the wet towels, the debris starts to rot and
stick to the towel. At this time they are more easily recognized and removed along with
a change of towel for the good seeds. Actually Acantholimon and most Asteraceae
germinate quickly at 70 so it is all over in a few weeks.

If the seed pods of Asclepias (milkweeds) are picked just before splitting, they
can be opened without fluffing up the silky parachutes attached to each seed.
Running the cylinder of seeds between the thumb and forefinger will allow (with luck)
the seed to be stripped of the silk to give seeds that are completely free of silk.
Seeds of Corydalis and Eranthis are shed from the capsule while the capsule is
still green. Close attention is needed to catch the seed before it falls. Seed of
Echinocereus (Cactaceae) is also shed from rather green fruit and need to be watched
closely. Seed of Violaceae has been some of the most difficult to collect: The capsule
is green one day with soft immature seed and only a few days later the capsule
suddenly splits and sends the seed flying. One solution is to wrap the immature
capsules in aluminum foil in a manner that catches the seed in the aluminum foil
envelope some days later. This was used with Adonis vernalis, Lathyrus vernus, and
many Violas.

Seeds of Buddleia and Weigelia are best collected by placing the stems lined
with the small seed capsules in a paper bag. They will open on drying and shed there
seeds whereas attempting to extract the seeds from the unopened capsules is most
difficult.

Seed of Phlox is frustrating to collect because just as it is ripe it rolls off the
recepticle, usually just as one is about to grab it. Also there are so few seeds per
recepticle. As a result seed of Phlox is usually scarce.

Seeds in fruits were washed as described in Chapter 8. After several days
washing and after most of the pulp had been mechanically removed, the seeds were
placed each day in a small strainer. By placing ones hands over the bottom of the
screen, the strainer was allowed to half fill and then drain by removing ones hand.
This was so efficient that 10-20 rinses per minute could be accomplished. The seeds
can then be transferred back to the soaking dish by inverting the strainer and using a
squirt bottle of the type used in chemistry laboratories.

One problem arises because nurserymen tend to propagate clones on the basis
of larger flowers, more fruit, and better foliage, and these are propagated vegetatively.
Gardeners also tend to plant a sort of Noah's Ark garden with one plant of each
species instead of self sowing colonies. In fact many gardeners look upon myriads of
self sown seedlings as disrupting their perfect picture. All of this leads to situations
that are not conducive to the production of viable seed in quantity. In our area
plantings of Pyracantha lalandi (firethorn) and Sorbus aucuparia (mountain ash) are
often composed entirely of a vegetatively propagated sterile clone. Abundant fruit is
set, which makes a great display but the fruits are seedless. Even more frustrating are
the trees of Photinia villosa on the. Penn State Campus that are laden with fruit filled
with normal size seed coats, but alas in many years all are empty shells. Beech trees
do the same and only occasionally is there a set of viable seeds. The only viable seed
of Pyracantha lalandi was found in a nearby junkyard where the birds had
undoubtedly roosted and sown the seed. Some species seem to be losing the ability
to set viable seed, clearly the case with Ceratostigma plumbaginoides and Franklinia.

The notion that all good seed sinks in water and bad seed floats is just not
always true. All the Iris setosa seed floats even after a month in water. In fact it starts
to germinate after a couple of weeks floating. Iris pseudoacorus seed floats for a few
days and then sinks. Large sample of seed of Cornus amomum were collected from
our own colonies and after thorough washing and cleaning, about half the seed floats
and half sinks. Both types gave about the same germination both in percent
germination and in the rates and other germination characters.

In some species much of the seed is infected with the eggs of certain insects
such as weevils. lpomoea leptophylla, the bush morning glory of the midwest U.S., is
particularly prone to this problem. One midwinter day I was working at my desk, and I
kept hearing a rattling and scraping sound. It was the weevils in a sack of the lpomoea
seeds. Certain beans have the same problem and one of our childhood playthings
were "Mexican Jumping Beans" Perhaps some carefully controlled heat treatment or
treatment with a toxic gas would kill the insects without damaging the seeds.

A disturbing problem is that due to extensive public spraying for gypsy moths
and oak leaf roller moths along with much private spraying of insecticides, the
population of moths and butterflies has been much depleted. I have not seen the
cocoon of one of the silk worm moths in years. In my childhood they were
commonplace. Of course the silkworm moths are not pollinators, but they reflect the
general decline in population of moths. Since the noctuid moths and the sphinx moths
are major pollinators of many flowers, the result is an absence of natural pollinators.
Fortunately the humming bird clear wing moth still is abundant and the hummingbird
itself is common and active, but the decline in population of insect pollinators has to be
a factor in reducing the seed set.

It was possible to obtain seeds from about 2500 species. Much of this seed was
personally collected. Our own plantings include woods, cliff, marsh, stream, and
spring and contain breeding colonies of over 400 species. The Penn State University
has a wide diversity of trees and shrubs on the campus, and the Borough of State
College employs a full time arborist and there are street plantings of interesting trees.
Most of the seeds were obtained from outside sources. The listing of these in
Chapter 23 may be helpful to gardeners. Foremost are the seed exchanges run by the
American flock Garden Society, the British Alpine Society, and the Scottish flock
Garden Club. These contain up to 6000 entries each year, and the Seedlist
Handbook (1) was written to sort out this plethora of names. Botanical Gardens were
important contributors with Khorag and Tbilisi in USSR, Beijing in China, and Denver
in US being particularly helpful. Many individuals took time to collect seeds and send
them to me and the list of over 100 are listed in Chapter 23. Finally there have been
botanical collecting expeditions that sent a wide variety of seeds. Notable were Jim
and Jenny Archibald in the Rockies and Turkey, Chris Chadwell in Kashmir and
Ladakh, Sonia and Jim Collins in Arizona, Jim Forrest in Australia and New Zealand,
Leon Glicenstein in California, Doris Goldman in Pennsylvania, Josef Halda in Turkey
and USSR, Gwen and Panayoti Kelaidis in the Rockies, VaclavPlestl in
Czechoslovakia, Ron Ratko in Washington, and Sally Walker in Southwestern US and
Mexico.

